\subsection{A stable reduction from finite categories to finite monoids}

The following construction preserves the local groups and shifts the
entire reduced cohomology, including its stable operations.  It is a
reduction for the general local-finiteness question; it does not assert
an unstable equivalence between the resulting classifying spaces.

\begin{proposition}\label{E:category-consolidation}
Let $\mathcal C$ be a nonempty finite category with $v$ objects and $a$
arrows.  There is an explicitly defined monoid $M_{\mathcal C}$ of order
$1+a+(v+1)^2$ with rectangular-band minimal ideal and a natural stable
equivalence
\[
 \widetilde{\Sigma^\infty BM_{\mathcal C}}
       \simeq \Sigma^2\Sigma^\infty(B\mathcal C)_+ .
\]
Every nontrivial maximal subgroup of $M_{\mathcal C}$ is a maximal
subgroup of an endomorphism monoid of $\mathcal C$, and conversely.
In particular, $M_{\mathcal C}$ is aperiodic if and only if all
endomorphism monoids of $\mathcal C$ are aperiodic.
\end{proposition}

\begin{proof}
Write $O=\operatorname{Ob}\mathcal C$ and $I=J=O\amalg\{*\}$.
Let $Q$ be the consolidation of $\mathcal C$: its elements are the
arrows of $\mathcal C$, a zero $0$, and a fresh identity $1$.
Multiplication of arrows is categorical composition when defined and
zero otherwise.  For an arrow $f:x\to y$, define maps
\[
 \sigma_f(x)=y,\qquad \sigma_f(z)=*\quad(z\ne x),
 \qquad
 \tau_f(y)=x,\qquad \tau_f(z)=*\quad(z\ne y).
\]
Both maps fix $*$.  Set $\sigma_1=\tau_1=\mathrm{id}$ and let the
zero act constantly at $*$.  Then
$\sigma_{fg}=\sigma_f\sigma_g$ and
$\tau_{fg}=\tau_g\tau_f$.

Take the disjoint union
\[
 M_{\mathcal C}=\{1\}\amalg\operatorname{Mor}\mathcal C
                       \amalg(I\times J).
\]
Give $I\times J$ its rectangular-band multiplication
$(i,j)(i',j')=(i,j')$.  An arrow $f$ acts by
$f(i,j)=(\sigma_f(i),j)$ and
$(i,j)f=(i,\tau_f(j))$.  The product of two arrows is their
categorical composite when it exists and $(*,*)$ otherwise.
The fresh $1$ is the identity.  The displayed action identities,
together with the fact that every arrow fixes $*$ on both sides,
verify associativity, including products which enter the band.
The band is the minimal ideal: multiplying any band element on the
left and right by band elements reaches every band element, and any
nonempty ideal contains a band element.  The quotient by this ideal
is precisely $Q$.

Put $A=\Sigma^\infty Q$, with the zero interpreted as the basepoint.
The stable band-bar formula of
Proposition~\ref{E:band-free-group-bar} identifies the left-hand side
with the pointed two-sided bar on the row and column free-group
spaces.  Here every nonzero arrow fixes the distinguished element
$*$, so its free-group map sends each basis generator to a basis
generator or to $1$; it does not introduce a conjugating word.
The canonical map from the wedge of basis circles to the classifying
space of the free group is natural for precisely these maps.
Consequently the two endpoint $A$-modules in that formula are
respectively $\Sigma U$ and $\Sigma V$, where
\[
 U=\bigvee_{x\in O}\mathbb S,
 \qquad V=\bigvee_{x\in O}\mathbb S
\]
carry the right and left partial-permutation actions above.  Thus
\[
 \widetilde{\Sigma^\infty BM_{\mathcal C}}
       \simeq \Sigma^2(U\mathbin{\wedge_A^{\mathbb L}}V).
\]
This use of the free-group model supplies the coherent module
identification; no choice of an equivariant splitting of the original
row augmentation is being assumed.

For each object $x$, the categorical identity $e_x$ is a literal
idempotent of $Q$.  Distinct $e_x,e_y$ are orthogonal.  The full
spectral category on these idempotent summands has morphism spectra
\[
 e_y A e_x=\Sigma^\infty\mathcal C(x,y)_+,
\]
with the given categorical composition.  Both endpoint modules are
supported on these summands, and each $e_x$-summand is $\mathbb S$.
The standard idempotent-corner bar comparison therefore identifies
$U\wedge_A^{\mathbb L}V$ with the two-sided constant bar of this
spectral category.  To see the comparison directly, split the
orthogonal idempotents in the stable module category and insert their
representable modules between consecutive factors.  Spectral Yoneda
contracts each adjacent representable pair; the resulting simplicial
term in degree $n$ is
\[
 \bigvee_{x_0\xrightarrow{f_1}x_1\longrightarrow\cdots
                   \xrightarrow{f_n}x_n}\mathbb S.
\]
Its face and degeneracy maps are exactly those of the nerve.  The
complementary fresh-identity summand acts trivially on both endpoints
and contributes no tensor factor.  Realization gives
$\Sigma^\infty(B\mathcal C)_+$, as required.  In particular, the
formal sum of the $e_x$ is used only as a stable idempotent, not as a
set map or as an additional element of $Q$.

Every idempotent of $M_{\mathcal C}$ is either the fresh identity, a
band element, or an idempotent endomorphism of $\mathcal C$.  The first
has trivial unit group, as no product of two nonidentity elements is
the fresh identity.  Band elements have trivial local groups.  At a
categorical idempotent $e$, an invertible element of the corner cannot
belong to the band, since a product involving the band cannot equal
$e$.  Its units are therefore exactly the units of the categorical
corner $e\operatorname{End}_{\mathcal C}(x)e$.  This proves the
assertion about local groups.
\end{proof}

\begin{corollary}\label{E:monoid-category-LF-equivalence}
Fix a prime $p$.  The assertion that the mod-$p$ cohomology of every
finite monoid whose maximal subgroups have order prime to $p$ is
locally finite over the Steenrod algebra is equivalent to the same
assertion for all finite categories whose local subgroups have order
prime to $p$.
\end{corollary}

\begin{proof}
One implication is immediate by viewing a monoid as a one-object
category.  For the other, apply
Proposition~\ref{E:category-consolidation}.  Stable suspension shifts
degree and commutes with every stable Steenrod operation; it neither
creates nor removes an infinite-dimensional operation orbit.  The
additional degree-zero class of $(B\mathcal C)_+$ causes no problem.
\end{proof}

\begin{corollary}\label{E:category-regular-nilpotent-scope}
Let $\mathcal C$ be a finite category whose pointed consolidation $Q$
has a nilpotent ideal $N$ with regular quotient $Q/N$.  At every prime
avoiding the orders of the maximal subgroups of $Q/N$, the mod-$p$
cohomology of $B\mathcal C$ is locally finite over the Steenrod algebra.
In particular $B\mathcal C$ cannot have the homotopy type of
$\Sigma\mathbb{CP}^{\infty}$.
\end{corollary}

\begin{proof}
The minimal ideal of $M_{\mathcal C}$ is a band and its pointed Rees
quotient is $Q$.  Apply
Theorem~\ref{E:nilpotent-ideal-regular-quotient} and the stable
equivalence above.  There is always a prime outside the finite set of
local-group divisors, whereas the suspended integral projective-space
generator has an infinite Steenrod orbit at every prime.
\end{proof}
